\documentclass[conference]{IEEEtran}
\IEEEoverridecommandlockouts

\usepackage{cite}
\usepackage{amsmath,amssymb,amsfonts}
\usepackage{graphicx}
\usepackage{textcomp}
\usepackage{xcolor}
\usepackage{booktabs}
\usepackage{multirow}
\usepackage{hyperref}
\usepackage{array}
\usepackage{balance}

\hypersetup{colorlinks=true,linkcolor=blue,citecolor=blue,urlcolor=blue}

\def\BibTeX{{\rm B\kern-.05em{\sc i\kern-.025em b}\kern-.08em
    T\kern-.1667em\lower.7ex\hbox{E}\kern-.125emX}}

\begin{document}

\title{Comparative Customer Segmentation Using Credit Card Behavioral Data:\\A Multi-Algorithm Clustering Analysis}

\author{
\IEEEauthorblockN{Om Tilekar}
\IEEEauthorblockA{
  \textit{MS in Data Science}\\
  \textit{University}\\
  om.tilekar@university.edu}
}

\maketitle

%---------------------------------------------------------------
\begin{abstract}
Customer segmentation is a fundamental task in financial services,
enabling targeted marketing, risk assessment, and personalized
retention strategies. This paper presents a comparative analysis of
four clustering algorithms---K-Means, Agglomerative Clustering,
DBSCAN, and Gaussian Mixture Models (GMM)---applied to a publicly
available credit card behavioral dataset of 8,950 customers described
by 17 features. A systematic preprocessing pipeline was developed
comprising median imputation, winsorization, log-transformation, and
StandardScaler normalization. Each algorithm was evaluated using
silhouette score, Davies-Bouldin index, and Calinski-Harabasz index.
K-Means ($k{=}4$) achieved the highest silhouette score (0.2099) and
Calinski-Harabasz index (2245.8), identifying four behaviorally
distinct segments: Installment Buyers, Cash-Dependent Users,
High-Value Transactors, and Low-Engagement Users. DBSCAN failed to
produce stable clusters across 24 hyperparameter configurations due
to the continuous, overlapping nature of financial behavioral data in
high-dimensional space. Results demonstrate that partition-based
methods outperform density-based approaches for credit card customer
segmentation, while Agglomerative Clustering provides strong
corroborating evidence for the four-cluster solution.
\end{abstract}

\begin{IEEEkeywords}
customer segmentation, K-Means, DBSCAN, GMM, agglomerative
clustering, credit card data, unsupervised learning, data mining
\end{IEEEkeywords}

%---------------------------------------------------------------
\section{Introduction}

Customer segmentation enables financial institutions to understand
behavioral heterogeneity within their customer base, supporting
data-driven decisions in credit risk management, product design, and
targeted marketing~\cite{ngai2009application}. Credit card transaction
data offers a rich behavioral signal---encompassing purchase patterns,
cash advance usage, payment regularity, and credit
utilization---making it well-suited for unsupervised clustering.

Clustering algorithms differ fundamentally in their assumptions.
Partition-based methods such as K-Means minimize within-cluster
variance assuming spherical clusters~\cite{macqueen1967methods}.
Hierarchical methods such as Agglomerative Clustering produce nested
groupings through successive merging~\cite{murtagh2012algorithms}.
Density-based methods such as DBSCAN identify arbitrarily shaped
clusters and handle noise without a pre-specified
$k$~\cite{ester1996density}. Probabilistic methods such as GMM
provide soft assignments via the Expectation-Maximization
algorithm~\cite{reynolds2009gaussian}. No single algorithm is
universally optimal, and comparative evaluation is essential for
robust segmentation.

This paper addresses: \textit{which clustering technique produces the
most meaningful, well-separated, and interpretable customer segments
from credit card usage data?} Our contributions are: (1) a systematic
preprocessing pipeline for financial behavioral data; (2) a
quantitative comparison of four clustering paradigms; and (3)
actionable segment profiles with direct business interpretability.

The remainder is organized as follows. Section~\ref{sec:related}
reviews related work. Section~\ref{sec:methodology} describes the
dataset and methodology. Section~\ref{sec:results} presents results.
Section~\ref{sec:discussion} discusses findings and limitations.
Section~\ref{sec:conclusion} concludes.

%---------------------------------------------------------------
\section{Related Work}
\label{sec:related}

Customer segmentation using clustering is well-studied in the
marketing and financial services
literature~\cite{ngai2009application}. K-Means has been the dominant
approach, consistently identifying 3--6 archetypes such as
transactors, revolvers, and dormant
customers~\cite{morissette2013k}. Its sensitivity to initialization
and assumption of spherical clusters motivates consideration of
alternatives.

Agglomerative Clustering provides a parameter-free structural view,
with dendrogram analysis guiding cluster count
selection~\cite{murtagh2012algorithms}. Ward linkage produces compact
clusters comparable to K-Means in financial
applications~\cite{nielsen2016introduction}. DBSCAN has been applied
successfully to spatial and network data but faces challenges in
high-dimensional tabular datasets where inter-point distances converge
and density contrasts
diminish~\cite{ester1996density,xu2015comprehensive}. GMM
accommodates elliptical cluster shapes and quantifies assignment
uncertainty, with BIC-based model selection as standard
practice~\cite{reynolds2009gaussian,fraley2002model}.

To our knowledge, no prior work systematically compares all four
paradigms---partitional, hierarchical, density-based, and
probabilistic---on the same credit card dataset with consistent
evaluation metrics and preprocessing, which motivates the present
study.

%---------------------------------------------------------------
\section{Methodology}
\label{sec:methodology}

\subsection{Dataset}

The Credit Card Customer Dataset (CCDATA) was obtained from
Kaggle~\cite{kaggle_ccdata}. It contains behavioral usage data for
8,950 active credit card holders over six months, described by 18
features including balance, purchase amounts, cash advance usage,
payment behavior, credit limit, and tenure. Two features contained
missing values: \texttt{MINIMUM\_PAYMENTS} ($n{=}313$, 3.5\%) and
\texttt{CREDIT\_LIMIT} ($n{=}1$). Approximately 47\% of customers
recorded zero total purchases and 59\% made no cash advances,
indicating a significant inactive customer population.

\subsection{Preprocessing Pipeline}

A six-stage pipeline was applied prior to all clustering experiments.

\textbf{Missing value imputation.} Both columns were imputed using
the column median, which is robust to the right-skewed distributions
characteristic of financial data. An assertion confirmed zero
remaining nulls after imputation.

\textbf{Column removal.} The identifier \texttt{CUST\_ID} was
dropped, yielding 17 behavioral features.

\textbf{Outlier mitigation.} Winsorization at the 1st and 99th
percentiles was applied using \texttt{scipy.stats.mstats.winsorize},
capping extremes without discarding rows and preserving
outlier-defined segments such as ultra-high spenders.

\textbf{Log-transformation.} Nine magnitude features with skewness
exceeding 2.0---including \texttt{BALANCE}, \texttt{PURCHASES},
\texttt{CASH\_ADVANCE}, \texttt{CREDIT\_LIMIT}, and related
transaction counts---were transformed using $\log(x{+}1)$. Frequency
and ratio features bounded in $[0,1]$ were excluded.
Fig.~\ref{fig:distributions} illustrates the skewness reduction
achieved.

\begin{figure}[t]
\centering
\includegraphics[width=\columnwidth]{fig3_distributions_fixed.png}
\caption{Feature distributions before (red) and after (green)
log-transformation for six representative features. Raw distributions
exhibit extreme right skew; log-transformation produces substantially
more symmetric distributions.}
\label{fig:distributions}
\end{figure}

\textbf{Feature scaling.} All 17 features were standardized using
\texttt{StandardScaler} to zero mean and unit variance. Without this
step, \texttt{CREDIT\_LIMIT} (range 0--30,000) would dominate
\texttt{TENURE} (range 6--12) in any distance computation.

\textbf{Dimensionality analysis.} PCA was applied to the scaled
matrix. Fig.~\ref{fig:pca} shows that six components explain 85\%
of cumulative variance and eight explain 92\%. All 17 features were
retained for clustering; 2-component PCA was used solely for
visualization.

\begin{figure}[t]
\centering
\includegraphics[width=\columnwidth]{fig4_pca_scree.png}
\caption{PCA cumulative explained variance scree plot. Six components
reach the 85\% threshold (red dashed); eight reach $\approx$92\%
(orange dashed). PCA was used for visualization only.}
\label{fig:pca}
\end{figure}

Fig.~\ref{fig:correlation} shows the Pearson correlation matrix of
scaled features. Strong associations are observed between
\texttt{PURCHASES\_TRX} and \texttt{PURCHASES\_FREQUENCY}
($r{=}0.91$), and \texttt{CASH\_ADVANCE\_FREQUENCY} and
\texttt{CASH\_ADVANCE\_TRX} ($r{=}0.92$), confirming coherent
behavioral dimensions.

\begin{figure}[t]
\centering
\includegraphics[width=\columnwidth]{fig2_correlation_heatmap.png}
\caption{Pearson correlation matrix of standardized features (lower
triangle). Strong correlations between purchase transaction count and
frequency ($r{=}0.91$) and between cash advance frequency and
transaction count ($r{=}0.92$) reflect coherent behavioral
dimensions.}
\label{fig:correlation}
\end{figure}

\subsection{Clustering Algorithms}

\textbf{K-Means} served as the baseline. The optimal $k$ was
determined jointly by WCSS elbow analysis and silhouette scores over
$k \in \{2,\ldots,12\}$, shown in Fig.~\ref{fig:elbow}. Although the
silhouette peaked at $k{=}2$ (0.263), this partition lacked
granularity. Cluster size analysis over $k \in \{3,\ldots,7\}$
identified $k{=}4$ as producing four well-balanced segments
(1,656--2,718 customers each). Final fitting used \texttt{k-means++}
initialization with $n\_init{=}20$.

\begin{figure}[t]
\centering
\includegraphics[width=\columnwidth]{fig_kmeans_elbow_silhouette.png}
\caption{K-Means cluster selection: WCSS elbow curve (left) and
silhouette score by $k$ (right). The smooth WCSS decline indicates a
continuous behavioral gradient. $k{=}4$ was selected based on cluster
balance, size uniformity, and interpretability despite the silhouette
peak at $k{=}2$.}
\label{fig:elbow}
\end{figure}

\textbf{Agglomerative Clustering} used ward linkage with $k{=}4$ for
direct comparison. Ward linkage minimizes within-cluster variance at
each merge, making it the most comparable hierarchical analogue to
K-Means. A truncated dendrogram on a 300-row sample confirmed the
cluster count.

\textbf{DBSCAN} was evaluated across 24 hyperparameter combinations
($\varepsilon \in \{0.5, 0.8, 1.0, 1.2, 1.5, 1.8, 2.0, 2.5, 3.0\}$,
$\text{min\_samples} \in \{5, 10, 15, 20, 33\}$). The best available
configuration ($\varepsilon{=}1.8$, $\text{min\_samples}{=}10$)
produced 3 clusters with 8.1\% noise and was retained for
completeness.

\textbf{GMM} with full covariance was fitted for $n \in \{2,\ldots,10\}$
components. Marginal BIC analysis identified $n{=}7$ as the point of
maximum curvature; BIC deteriorated at $n{=}8$ before recovering
irregularly, indicating additional components model noise rather than
genuine structure.

\subsection{Evaluation Metrics}

Three internal validation metrics were computed for all models:
\textbf{Silhouette Score}~\cite{rousseeuw1987silhouettes} (cohesion
and separation; $\uparrow$ better, range $[-1,1]$);
\textbf{Davies-Bouldin Index}~\cite{davies1979cluster} (average
cluster similarity; $\downarrow$ better);
\textbf{Calinski-Harabasz Index}~\cite{calinski1974dendrite}
(between-to-within dispersion ratio; $\uparrow$ better). Noise
points ($\text{label}{=}{-1}$) were excluded from all DBSCAN metric
computations.

%---------------------------------------------------------------
\section{Results}
\label{sec:results}

\subsection{Algorithm Comparison}

Table~\ref{tab:comparison} presents evaluation metrics for all four
algorithms. K-Means achieved the highest silhouette score (0.2099)
and Calinski-Harabasz index (2245.8), winning two of three metrics.
Agglomerative Clustering corroborated the four-cluster solution
independently (silhouette: 0.1774, CH: 1870.3). DBSCAN's
Davies-Bouldin index (0.9681) reflects evaluation on only 91.9\%
non-noise points and is not directly comparable. GMM ($n{=}7$)
achieved finer granularity at the cost of compactness (silhouette:
0.1644).

\begin{table}[t]
\centering
\caption{Clustering Algorithm Comparison}
\label{tab:comparison}
\resizebox{\columnwidth}{!}{%
\begin{tabular}{@{}lcccc@{}}
\toprule
Model         & $k$ & Silhouette$\uparrow$ & Davies-Bouldin$\downarrow$ & Calinski-Harabasz$\uparrow$ \\
\midrule
K-Means       & 4   & \textbf{0.2099}      & 1.6708                     & \textbf{2245.8}              \\
Agglomerative & 4   & 0.1774               & 1.7218                     & 1870.3                       \\
DBSCAN$^{*}$  & 3   & 0.1160               & \textbf{0.9681}            & 14.5                         \\
GMM           & 7   & 0.1644               & 2.3007                     & 1334.6                       \\
\bottomrule
\multicolumn{5}{l}{\footnotesize $^{*}$Metrics exclude 8.1\% noise points; DB index not directly comparable.}
\end{tabular}%
}
\end{table}

Fig.~\ref{fig:comparison} presents PCA 2D projections of all four
solutions. K-Means and Agglomerative produce nearly identical spatial
boundaries, strongly validating the four-cluster structure. DBSCAN
assigns the majority to a single dense region. GMM subdivides K-Means
regions into finer components.

\begin{figure*}[t]
\centering
\includegraphics[width=\textwidth]{fig_all_models_comparison.png}
\caption{PCA 2D projections of all four clustering solutions. K-Means
and Agglomerative Clustering produce nearly identical spatial
boundaries, validating the four-cluster partition. DBSCAN assigns
most customers to a single cluster with scattered noise points. GMM
subdivides the K-Means regions into seven finer components,
particularly splitting high-purchase and high-balance zones.}
\label{fig:comparison}
\end{figure*}

\subsection{Customer Segment Profiles}

Table~\ref{tab:cluster_sizes} shows the four K-Means segment sizes.
All are substantially populated (18.5\%--30.4\%), confirming $k{=}4$
does not produce degenerate micro-clusters. Fig.~\ref{fig:sizes}
visualizes the distribution.

\begin{table}[t]
\centering
\caption{K-Means Cluster Sizes ($k{=}4$)}
\label{tab:cluster_sizes}
\begin{tabular}{@{}llcc@{}}
\toprule
Cluster & Segment Name           & Count          & Share   \\
\midrule
C0      & Installment Buyers     & 2,141          & 23.9\%  \\
C1      & Cash-Dependent Users   & 2,718          & 30.4\%  \\
C2      & High-Value Transactors & 2,435          & 27.2\%  \\
C3      & Low-Engagement Users   & 1,656          & 18.5\%  \\
\midrule
        & \textbf{Total}         & \textbf{8,950} & \textbf{100\%} \\
\bottomrule
\end{tabular}
\end{table}

\begin{figure}[t]
\centering
\includegraphics[width=\columnwidth]{fig_cluster_sizes.png}
\caption{Customer segment distribution for K-Means ($k{=}4$). All
four segments contain 1,656--2,718 customers, confirming no
degenerate micro-clusters.}
\label{fig:sizes}
\end{figure}

Table~\ref{tab:profiles} reports key centroid values in original
scale. Fig.~\ref{fig:centroid} shows the full standardized centroid
heatmap; Fig.~\ref{fig:radar} shows radar profiles across eight
behavioral dimensions.

\begin{table}[t]
\centering
\caption{K-Means Cluster Centroids (Selected Features, Original Scale)}
\label{tab:profiles}
\resizebox{\columnwidth}{!}{%
\begin{tabular}{@{}lcccc@{}}
\toprule
Feature              & C0 Installment & C1 Cash-Dep. & C2 High-Val. & C3 Low-Eng. \\
\midrule
BALANCE              & 4.97    & 7.39    & 7.16    & 4.21   \\
PURCHASES            & 6.16    & 1.57    & 7.45    & 4.97   \\
CASH\_ADVANCE        & 0.72    & 6.99    & 3.10    & 0.96   \\
INSTALLMENTS\_PURCH. & 6.04    & 0.48    & 5.31    & 1.68   \\
PURCHASES\_FREQ.     & 0.82    & 0.07    & 0.85    & 0.23   \\
CREDIT\_LIMIT        & 7.73    & 8.09    & 8.56    & 7.90   \\
PAYMENTS             & 802.83  & 1705.62 & 2895.94 & 821.21 \\
PRC\_FULL\_PAYMENT   & 0.29    & 0.03    & 0.17    & 0.15   \\
\bottomrule
\end{tabular}%
}
\end{table}

\begin{figure}[t]
\centering
\includegraphics[width=\columnwidth]{fig_centroid_heatmap.png}
\caption{K-Means cluster centroids as standardized feature values.
Blue = above-average; red = below-average for that cluster. Each row
exhibits a distinct pattern confirming behavioral differentiation
across all four segments.}
\label{fig:centroid}
\end{figure}

\begin{figure}[t]
\centering
\includegraphics[width=\columnwidth]{fig_radar_charts.png}
\caption{Radar chart profiles for the four K-Means segments across
eight normalized behavioral dimensions. C2 (High-Value Transactors)
fills the largest polygon; C1 (Cash-Dependent Users) shows a
cross-shaped spike on cash advance; C0 (Installment Buyers) shows a
narrow purchase-frequency wedge; C3 (Low-Engagement Users) produces
the smallest polygon.}
\label{fig:radar}
\end{figure}

\textbf{C0---Installment Buyers (23.9\%).} High installment purchase
activity (6.04), high purchase frequency (0.82), highest full-payment
rate (0.29), and minimal cash advance (0.72). These disciplined,
low-risk customers buy regularly through structured payment plans.

\textbf{C1---Cash-Dependent Users (30.4\%).} The largest segment
exhibits high cash advance amounts (6.99), cash advance frequency
(0.30), near-zero purchases (1.57, frequency: 0.07), and the lowest
full-payment rate (0.03). These customers use their card primarily as
a liquidity instrument, signaling potential financial stress and
elevated credit risk.

\textbf{C2---High-Value Transactors (27.2\%).} The premium segment
records the highest purchases (7.45), one-off purchases (6.64),
credit limit (8.56), and total payments (2,895.94). The radar polygon
covers the largest area across all dimensions, confirming these as the
most financially active and institutionally valuable customers.

\textbf{C3---Low-Engagement Users (18.5\%).} Uniformly low values
across all dimensions---below-average purchases (4.97), low frequency
(0.23), minimal cash advances, and the lowest balance frequency
(0.60). The radar polygon is the most compressed, confirming minimal
card engagement and representing a significant activation opportunity.

Fig.~\ref{fig:parallel} corroborates these profiles via parallel
coordinates, with clearly separated trajectory bundles especially at
\texttt{CASH\_ADVANCE} and \texttt{PURCHASES\_FREQUENCY}.

\begin{figure}[t]
\centering
\includegraphics[width=\columnwidth]{fig_parallel_coordinates.png}
\caption{Parallel coordinates of all four customer segments
(standardized scale, 300-customer sample per segment). Separated
trajectory bundles at \texttt{CASH\_ADVANCE} and
\texttt{PURCHASES\_FREQUENCY} corroborate the behavioral
distinctiveness of each segment.}
\label{fig:parallel}
\end{figure}

%---------------------------------------------------------------
\section{Discussion}
\label{sec:discussion}

\subsection{Algorithm Suitability for Financial Behavioral Data}

The convergence of K-Means and Agglomerative Clustering on the same
four-cluster spatial partition---despite fundamentally different
algorithmic approaches---provides strong cross-validation that the
four-segment solution captures genuine data structure rather than an
algorithmic artifact. This constitutes a methodological strength of
the comparative design.

DBSCAN's failure across 24 hyperparameter combinations is an
informative comparative finding. At low $\varepsilon$, up to 82.8\%
of customers were classified as noise. At high $\varepsilon$, all
customers collapsed into a single cluster. This behavior is consistent
with the curse of dimensionality in 17-dimensional continuous data,
where inter-point distances converge and eliminate the density
contrasts DBSCAN requires~\cite{xu2015comprehensive}. This suggests
that partition-based and hierarchical methods are better suited to
financial behavioral segmentation than density-based approaches.

GMM's $n{=}7$ components subdivide the K-Means partition into finer
groupings, capturing soft boundaries that K-Means cannot represent.
However, the reduced silhouette score (0.1644 vs.\ 0.2099) indicates
this additional complexity trades cluster compactness for
granularity. BIC deterioration at $n{=}8$ confirms $n{=}7$ as the
principled stopping point.

\subsection{Business Implications}

The four segments carry direct actionable interpretations.
Cash-Dependent Users (30.4\%) represent the highest credit risk;
proactive outreach, financial counseling, or credit limit adjustment
may reduce default exposure. High-Value Transactors (27.2\%) are the
most profitable segment and should be prioritized for premium products
and loyalty programs. Installment Buyers (23.9\%) are low-risk and
may respond to installment plan expansion and balance transfer offers.
Low-Engagement Users (18.5\%) represent a significant activation
opportunity via targeted incentive campaigns.

\subsection{Limitations}

Several limitations apply. First, the dataset covers a single
six-month behavioral window; longitudinal data would enable customer
migration analysis across segments over time. Second, all validation
metrics are internal; no ground-truth labels exist for external
validation, which is inherent to unsupervised learning. Third,
zero-inflation in purchase and cash advance features (47\% and 59\%
zero values) may conflate behaviorally distinct inactive customers
into the single low-engagement segment. Fourth, PCA applied before
DBSCAN may improve density separation and constitutes a promising
direction for future work. Fifth, the absence of sociodemographic
features limits segment profiling depth.

\subsection{Future Work}

Future directions include: (1) self-organizing maps or deep
clustering methods to handle continuous behavioral gradients;
(2) temporal modeling of customer segment transitions over time;
(3) semi-supervised approaches incorporating partial label information
from customer service outcomes; and (4) ensemble clustering to produce
consensus partitions more robust than any individual algorithm.

%---------------------------------------------------------------
\section{Conclusion}
\label{sec:conclusion}

This paper presented a systematic comparative evaluation of four
clustering algorithms for credit card customer segmentation.
Following a principled preprocessing pipeline addressing missing
values, outliers, skewness, and scale, all methods were applied to
8,950 customers across 17 behavioral features.

K-Means ($k{=}4$) was the best-performing algorithm (silhouette:
0.2099, CH: 2245.8), producing four interpretable segments with clear
business implications. Its spatial convergence with Agglomerative
Clustering validates the robustness of the four-cluster solution.
DBSCAN's instability across 24 configurations highlights the
limitations of density-based methods on high-dimensional financial
behavioral data. GMM offered probabilistic granularity at the cost of
cluster compactness.

The four identified segments---Installment Buyers, Cash-Dependent
Users, High-Value Transactors, and Low-Engagement Users---correspond
to actionable archetypes that can directly inform credit risk
management, product design, and targeted marketing in financial
services. Code and notebooks are available at the project GitHub
repository.

%---------------------------------------------------------------
\balance
\begin{thebibliography}{00}

\bibitem{ngai2009application}
E.\,W.\,T. Ngai, L. Xiu, and D.\,C.\,K. Chau,
``Application of data mining techniques in customer relationship
management: A literature review and classification,''
\textit{Expert Systems with Applications},
vol.~36, no.~2, pp.~2592--2602, 2009.

\bibitem{macqueen1967methods}
J. MacQueen,
``Some methods for classification and analysis of multivariate
observations,''
in \textit{Proc. 5th Berkeley Symp. Mathematical Statistics and
Probability}, vol.~1, pp.~281--297, 1967.

\bibitem{murtagh2012algorithms}
F. Murtagh and P. Contreras,
``Algorithms for hierarchical clustering: An overview,''
\textit{WIREs Data Mining and Knowledge Discovery},
vol.~2, no.~1, pp.~86--97, 2012.

\bibitem{ester1996density}
M. Ester, H.-P. Kriegel, J. Sander, and X. Xu,
``A density-based algorithm for discovering clusters in large
spatial databases with noise,''
in \textit{Proc. 2nd Int. Conf. Knowledge Discovery and Data Mining
(KDD)}, pp.~226--231, 1996.

\bibitem{reynolds2009gaussian}
D.\,A. Reynolds,
``Gaussian mixture models,''
\textit{Encyclopedia of Biometrics}, pp.~659--663, 2009.

\bibitem{kaggle_ccdata}
A. Bhasin,
``Credit Card Dataset for Clustering,''
Kaggle, 2019. [Online]. Available:
\url{https://www.kaggle.com/arjunbhasin2013/ccdata}

\bibitem{rousseeuw1987silhouettes}
P.\,J. Rousseeuw,
``Silhouettes: A graphical aid to the interpretation and validation
of cluster analysis,''
\textit{Journal of Computational and Applied Mathematics},
vol.~20, pp.~53--65, 1987.

\bibitem{davies1979cluster}
D.\,L. Davies and D.\,W. Bouldin,
``A cluster separation measure,''
\textit{IEEE Trans. Pattern Analysis and Machine Intelligence},
vol.~PAMI-1, no.~2, pp.~224--227, 1979.

\bibitem{calinski1974dendrite}
T. Cali\'{n}ski and J. Harabasz,
``A dendrite method for cluster analysis,''
\textit{Communications in Statistics},
vol.~3, no.~1, pp.~1--27, 1974.

\bibitem{morissette2013k}
L. Morissette and S. Chartier,
``The k-means clustering technique: General considerations and
implementation in Mathematica,''
\textit{Tutorials in Quantitative Methods for Psychology},
vol.~9, no.~1, pp.~15--24, 2013.

\bibitem{nielsen2016introduction}
F. Nielsen,
\textit{Introduction to HPC with MPI for Data Science}.
Springer, 2016.

\bibitem{xu2015comprehensive}
D. Xu and Y. Tian,
``A comprehensive survey of clustering algorithms,''
\textit{Annals of Data Science},
vol.~2, no.~2, pp.~165--193, 2015.

\bibitem{fraley2002model}
C. Fraley and A.\,E. Raftery,
``Model-based clustering, discriminant analysis, and density
estimation,''
\textit{Journal of the American Statistical Association},
vol.~97, no.~458, pp.~611--631, 2002.

\end{thebibliography}

\end{document}
